\section{Síntesis y ampliación}

\subsection{Repaso de los puntos clave}

Esta clase se centró en \textbf{Score-Based Models}, estableciendo sistemáticamente los siguientes conceptos fundamentales:

\begin{enumerate}
    \item \textbf{Score Function}: definido como $\nabla_x \log p(x)$, es un campo vectorial que apunta en la dirección de aumento de probabilidad y codifica completamente la información de la distribución de probabilidad.

    \item \textbf{Predicción de ruido como aprendizaje del Score}: el predictor de ruido $\epsilon_\theta$ en DDPM difiere del score function únicamente por un factor de escala, revelando el profundo significado físico de la predicción de ruido.

    \item \textbf{Fórmula de Tweedie}: proporciona una fórmula directa desde el score function hasta la estimación de la media posterior (denoising): $\mathbb{E}[\mu \mid x] = x + \sigma^2 \nabla_x \log p(x)$.

    \item \textbf{Dinámicas de Langevin}: un método de muestreo que depende únicamente del score function, convergiendo a la distribución objetivo mediante un proceso iterativo de ``ascenso por gradiente + inyección de ruido''.

    \item \textbf{Score Matching con desruido}: transforma el objetivo no computable de score matching en un objetivo computable de desruido; su proceso de entrenamiento es totalmente equivalente al de DDPM.

    \item \textbf{Ruido multiescala y estrategia de recocido}: resuelve el problema de inexactitud del score en regiones de baja densidad mediante múltiples niveles de ruido del score function y dinámicas de Langevin con recocido, lo cual es equivalente al proceso inverso de desruido de DDPM.
\end{enumerate}

\subsection{Lecturas adicionales}

\begin{itemize}
    \item \textbf{Song \& Ermon (2019)}: \textit{Generative Modeling by Estimating Gradients of the Data Distribution} -- propone NCSN, integrando por primera vez score matching y dinámicas de Langevin con recocido para la generación de imágenes.
    \item \textbf{Ho et al. (2020)}: \textit{Denoising Diffusion Probabilistic Models (DDPM)} -- deriva el objetivo de entrenamiento de los modelos de difusión desde una perspectiva de inferencia variacional.
    \item \textbf{Song et al. (2021)}: \textit{Score-Based Generative Modeling through Stochastic Differential Equations} -- propone el marco Score SDE, unificando DDPM y NCSN como procesos continuos en el tiempo.
    \item \textbf{Vincent (2011)}: \textit{A Connection Between Score Matching and Denoising Autoencoders} -- establece por primera vez los fundamentos teóricos del score matching con desruido.
    \item \textbf{Efron (2011)}: \textit{Tweedie's Formula and Selection Bias} -- un tratamiento clásico de la fórmula de Tweedie en estadística.
\end{itemize}

%% --- Fin del contenido --- %%

\end{document}
